% TOP_PROBLEM: {"id": 30002559, "title": "Spectrahedrality of Hyperbolicity Cones", "category_id": 6, "category": {"id": 6, "name": "geometry", "display_name": "Geometry", "description": "Euclidean and non-Euclidean geometry, geometric structures.", "slug": "geometry", "order_index": 6, "created_at": "2026-07-31T15:26:25.671Z"}, "status": "open"}
% FIRST_POSTED: 2026-10-06
% ATTEMPT_STATUS: solved
% OPENAI_CLAIM: {"schema_version":1,"source_repo":"https://github.com/openai/math","source_commit":"adc7f1241b42e322a6451854ab7e4b4c146bf78a","release_date":"2026-10-06","match":"stronger","resolution":"disproved","independently_reviewed":false,"second_pass":"agreed","families":[{"family":"095","title":"Hyperbolicity cones without semidefinite lifts","subject":"Convex and metric geometry","manuscripts":[{"title":"A nonspectrahedral hyperbolicity cone","pdf_path":"preprints/A-Nonspectrahedral-Hyperbolicity-Cone-September-24-2026/nonspectrahedral-hyperbolicity-cone.pdf","readme_path":"preprints/A-Nonspectrahedral-Hyperbolicity-Cone-September-24-2026/README.md","date":"2026-09-24","theorem_ref":"thm:main (sections/01-introduction.tex:63)","statement_sources":[{"path":"preprints/A-Nonspectrahedral-Hyperbolicity-Cone-September-24-2026/build/sections/01-introduction.tex","line":63,"label":"thm:main"}]},{"title":"An Exact Semidefinite Lift of a Nonspectrahedral Hyperbolicity Cone","pdf_path":"preprints/An-Exact-Semidefinite-Lift-of-a-Nonspectrahedral-Hyperbolicity-Cone-October-5-2026/Exact-Semidefinite-Lift-of-a-Nonspectrahedral-Hyperbolicity-Cone.pdf","readme_path":"preprints/An-Exact-Semidefinite-Lift-of-a-Nonspectrahedral-Hyperbolicity-Cone-October-5-2026/README.md","date":"2026-10-05","theorem_ref":"thm:main (sections/introduction.tex:68)","statement_sources":[{"path":"preprints/An-Exact-Semidefinite-Lift-of-a-Nonspectrahedral-Hyperbolicity-Cone-October-5-2026/build/sections/introduction.tex","line":68,"label":"thm:main"}]},{"title":"Hyperbolicity Cones Without Semidefinite Lifts","pdf_path":"preprints/Hyperbolicity-Cones-Without-Semidefinite-Lifts-October-5-2026/nonliftable-hyperbolicity.pdf","readme_path":"preprints/Hyperbolicity-Cones-Without-Semidefinite-Lifts-October-5-2026/README.md","date":"2026-10-05","theorem_ref":"thm:main (source/01-introduction.tex:33)","statement_sources":[{"path":"preprints/Hyperbolicity-Cones-Without-Semidefinite-Lifts-October-5-2026/build/source/01-introduction.tex","line":33,"label":"thm:main"}]}],"lean":{"doc_path":"lean/docs/095.md","comparators":["lean/ComparatorChallenges/HyperbolicCones.json","lean/ComparatorChallenges/HyperbolicCones.lean"],"declarations":["OAI.Paper256.main_result"],"covers_main_theorem":false},"reasoning_trace":null}]}
\documentclass[11pt]{article}
\usepackage[margin=25mm]{geometry}
\usepackage{fontspec}
\setmainfont{Latin Modern Roman}
\usepackage{amsmath,amssymb,mathtools}
\usepackage{unicode-math}
\setmathfont{Latin Modern Math}
\usepackage{enumitem}
\usepackage{xurl,hyperref,bookmark}
\hypersetup{colorlinks=true,urlcolor=blue}
\setcounter{secnumdepth}{0}
\setlength{\parindent}{0pt}
\setlength{\parskip}{5pt}
\setlength{\emergencystretch}{3em}
\begin{document}
\section*{Spectrahedrality of Hyperbolicity Cones}

\subsection{OpenAI's claim}
OpenAI presents the following claims relevant to this question, with manuscript numbers following paper order in Sources. In manuscript 1 (family 095; 2026-09-24), the asserted result is that an explicit degree-sixteen polynomial in twenty-three real variables is hyperbolic but its cone has no finite homogeneous real-symmetric matrix-inequality representation in the original coordinates; In manuscript 2 (family 095; 2026-10-05), the asserted result is that the explicit twenty-three-variable nonspectrahedral cone nonetheless has an exact lift with a size-$100$ pencil and $307$ auxiliary variables, including singular boundary points; In manuscript 3 (family 095; 2026-10-05), the asserted result is that some hyperbolicity cones remain outside the class of projections of finite real semidefinite systems, whatever real coefficients and auxiliary variables are used.

\subsection{Mathematical statement of the claimed result}
\paragraph{Manuscript 1.}
Use cyclic indices on $y\in\mathbb R^3$ and set $Q(y)_{ii}=y_i^2+y_{i+1}^2$, $Q(y)_{ij}=-y_i y_j$ for $i\ne j$. For real symmetric $4\times4$ matrices, define
\[
 \Phi_y\begin{pmatrix}a&r^{\mathsf T}\\r&T\end{pmatrix}
 =\begin{pmatrix}\operatorname{tr}(Q(y)T)&-r^{\mathsf T}Q(y)\\-Q(y)r&aQ(y)\end{pmatrix},
 \quad p(X,Z,y)=\det((\det X)Z-\Phi_y(\operatorname{adj}X)).
\]
The claimed polynomial has degree $20$, satisfies $p(I_4,I_4,0)=1$, and is hyperbolic in that direction. Its closed hyperbolicity cone $K$ admits no direct linear matrix inequality: for every finite $N\ge1$ and every real linear $L$ from the $23$-dimensional ambient space to real symmetric $N\times N$ matrices,
\[
 K\ne\{(X,Z,y):L(X,Z,y)\succeq0\}.
\]
This excludes direct spectrahedrality, rather than merely a particular chosen pencil.

\paragraph{Manuscript 2.}
For cyclic $i\in\{1,2,3\}$ set $Q(y)_{ii}=y_i^2+y_{i+1}^2$, $Q(y)_{ij}=-y_i y_j$ otherwise, and
\[
 \Phi_y\begin{pmatrix}a&r^{\mathsf T}\\r&T\end{pmatrix}
 =\begin{pmatrix}\operatorname{tr}(Q(y)T)&-r^{\mathsf T}Q(y)\\-Q(y)r&aQ(y)\end{pmatrix},
 \quad p(X,Z,y)=\det((\det X)Z-\Phi_y(\operatorname{adj}X)).
\]
Let $K$ be the closed hyperbolicity cone of $p$ in direction $(I_4,I_4,0)$ on the real $23$-dimensional space of symmetric $X,Z$ and $y\in\mathbb R^3$. The claim gives fixed real symmetric $100\times100$ matrices $A_1,\ldots,A_{23},B_1,\ldots,B_{307}$ with
\[
 x\in K\ \Longleftrightarrow\ \exists v\in\mathbb R^{307}:\quad
 \sum_{\nu=1}^{23}x_{\nu} A_{\nu}+\sum_{j=1}^{307}v_jB_j\succeq0.
\]
Coordinates list upper-triangular entries of $X$ then $Z$, followed by $y$. Singular boundary points are included. Auxiliary variables make this a lift; it does not assert direct spectrahedrality.

\paragraph{Manuscript 3.}
The existence claim is a real homogeneous polynomial $p$ and a direction $e$ with $p(e)>0$ such that every univariate $t\mapsto p(te-x)$ is real-rooted, but no finite semidefinite lift represents the resulting closed hyperbolicity cone $K$. Precisely,
\[
 \nexists N,r<\infty,\ A_i,B_j\in\operatorname{Sym}_N(\mathbb R):\quad
 K=\{x:\exists y\in\mathbb R^r,\ A_0+\sum_ix_iA_i+\sum_jy_jB_j\succeq0\}.
\]
This excludes representations even with auxiliary variables, and is stronger than excluding a direct linear matrix inequality.

\subsection{Outline of the argument}
For manuscript 1 (family 095; 2026-09-24), a nonnegative Choi--Lam biquadratic form with no dominated bilinear square is inserted into a determinant construction. Rank-one cone limits force a norm identity under a hypothetical pencil representation, and the resulting tangent obstruction contradicts the scalar form's extremality.

For manuscript 2 (family 095; 2026-10-05), the three-variable parameter is covered by rotating two-dimensional planes. Each angular slice has a trigonometric Gram factorization and an exact pencil representation; a finite matching identity combines these slices into the full homogeneous lift.

For manuscript 3 (family 095; 2026-10-05), a positive semidefinite Hankel seed with controlled multiplication kernel produces finitely many inverse-limit directions. Separation from Hankel images yields a polynomial patch whose moment deformation contradicts analytic data that any semidefinite lift would require.

\subsection{Scope relative to this problem}
The negative manuscript claims a hyperbolicity cone with no finite semidefinite lift. Since every spectrahedral cone has such a lift, this is a counterexample to the universal spectrahedrality question and asserts a stronger obstruction. A separate linked construction gives an exact semidefinite lift for its particular nonspectrahedral cone; a lift with auxiliary variables does not make that cone itself spectrahedral.

\subsection{Formal verification}
For manuscript 1 (family 095; 2026-09-24), the documented Lean scope is a degree-twenty hyperbolic polynomial in twenty-three variables whose cone has no pencil representation; this is not a formal transcription of the paper's degree-sixteen formula and does not assert nonexistence of lifts.

This coverage does not establish a formal proof of the complete problem as stated here.

No Lean formalization is listed at this snapshot for manuscript 2 (family 095; 2026-10-05), manuscript 3 (family 095; 2026-10-05).

The linked family Comparator checks are \texttt{HyperbolicCones.lean}.

Their named declarations include \texttt{OAI.Paper256.main\_result}.

This site has not executed the Lean code. Described statement coverage is not an independent proof check.

\subsection{Status caveat}
The manuscripts were produced by an internal OpenAI model. OpenAI cautions that results without formalization can contain mistakes. This site has not independently checked the proofs; the label records the scope of OpenAI's claim.

\subsection{Sources}
\begin{itemize}
\item OpenAI, \emph{A nonspectrahedral hyperbolicity cone}, 2026-09-24; citation key \texttt{OAI:A-Nonspectrahedral-Hyperbolicity-Cone-September-24-2026}. \href{https://github.com/openai/math/blob/adc7f1241b42e322a6451854ab7e4b4c146bf78a/preprints/A-Nonspectrahedral-Hyperbolicity-Cone-September-24-2026/nonspectrahedral-hyperbolicity-cone.pdf}{PDF (GitHub)}; \href{https://github.com/openai/math/raw/adc7f1241b42e322a6451854ab7e4b4c146bf78a/preprints/A-Nonspectrahedral-Hyperbolicity-Cone-September-24-2026/nonspectrahedral-hyperbolicity-cone.pdf}{PDF (direct)}; \href{https://github.com/openai/math/blob/adc7f1241b42e322a6451854ab7e4b4c146bf78a/preprints/A-Nonspectrahedral-Hyperbolicity-Cone-September-24-2026/README.md}{Manuscript page and citation}; \href{https://github.com/openai/math/blob/main/preprints/A-Nonspectrahedral-Hyperbolicity-Cone-September-24-2026/nonspectrahedral-hyperbolicity-cone.pdf}{Latest version}; Statement source: \href{https://github.com/openai/math/blob/adc7f1241b42e322a6451854ab7e4b4c146bf78a/preprints/A-Nonspectrahedral-Hyperbolicity-Cone-September-24-2026/build/sections/01-introduction.tex\#L63}{thm:main}.
\item OpenAI, \emph{An Exact Semidefinite Lift of a Nonspectrahedral Hyperbolicity Cone}, 2026-10-05; citation key \texttt{OAI:An-Exact-Semidefinite-Lift-of-a-Nonspectrahedral-Hyperbolicity-Cone-October-5-2026}. \href{https://github.com/openai/math/blob/adc7f1241b42e322a6451854ab7e4b4c146bf78a/preprints/An-Exact-Semidefinite-Lift-of-a-Nonspectrahedral-Hyperbolicity-Cone-October-5-2026/Exact-Semidefinite-Lift-of-a-Nonspectrahedral-Hyperbolicity-Cone.pdf}{PDF (GitHub)}; \href{https://github.com/openai/math/raw/adc7f1241b42e322a6451854ab7e4b4c146bf78a/preprints/An-Exact-Semidefinite-Lift-of-a-Nonspectrahedral-Hyperbolicity-Cone-October-5-2026/Exact-Semidefinite-Lift-of-a-Nonspectrahedral-Hyperbolicity-Cone.pdf}{PDF (direct)}; \href{https://github.com/openai/math/blob/adc7f1241b42e322a6451854ab7e4b4c146bf78a/preprints/An-Exact-Semidefinite-Lift-of-a-Nonspectrahedral-Hyperbolicity-Cone-October-5-2026/README.md}{Manuscript page and citation}; \href{https://github.com/openai/math/blob/main/preprints/An-Exact-Semidefinite-Lift-of-a-Nonspectrahedral-Hyperbolicity-Cone-October-5-2026/Exact-Semidefinite-Lift-of-a-Nonspectrahedral-Hyperbolicity-Cone.pdf}{Latest version}; Statement source: \href{https://github.com/openai/math/blob/adc7f1241b42e322a6451854ab7e4b4c146bf78a/preprints/An-Exact-Semidefinite-Lift-of-a-Nonspectrahedral-Hyperbolicity-Cone-October-5-2026/build/sections/introduction.tex\#L68}{thm:main}.
\item OpenAI, \emph{Hyperbolicity Cones Without Semidefinite Lifts}, 2026-10-05; citation key \texttt{OAI:Hyperbolicity-Cones-Without-Semidefinite-Lifts-October-5-2026}. \href{https://github.com/openai/math/blob/adc7f1241b42e322a6451854ab7e4b4c146bf78a/preprints/Hyperbolicity-Cones-Without-Semidefinite-Lifts-October-5-2026/nonliftable-hyperbolicity.pdf}{PDF (GitHub)}; \href{https://github.com/openai/math/raw/adc7f1241b42e322a6451854ab7e4b4c146bf78a/preprints/Hyperbolicity-Cones-Without-Semidefinite-Lifts-October-5-2026/nonliftable-hyperbolicity.pdf}{PDF (direct)}; \href{https://github.com/openai/math/blob/adc7f1241b42e322a6451854ab7e4b4c146bf78a/preprints/Hyperbolicity-Cones-Without-Semidefinite-Lifts-October-5-2026/README.md}{Manuscript page and citation}; \href{https://github.com/openai/math/blob/main/preprints/Hyperbolicity-Cones-Without-Semidefinite-Lifts-October-5-2026/nonliftable-hyperbolicity.pdf}{Latest version}; Statement source: \href{https://github.com/openai/math/blob/adc7f1241b42e322a6451854ab7e4b4c146bf78a/preprints/Hyperbolicity-Cones-Without-Semidefinite-Lifts-October-5-2026/build/source/01-introduction.tex\#L33}{thm:main}.
\item Family 095: \href{https://github.com/openai/math/blob/adc7f1241b42e322a6451854ab7e4b4c146bf78a/lean/docs/095.md}{Lean scope}.
\item \href{https://github.com/openai/math/blob/adc7f1241b42e322a6451854ab7e4b4c146bf78a/lean/ComparatorChallenges/HyperbolicCones.json}{Comparator HyperbolicCones.json}.
\item \href{https://github.com/openai/math/blob/adc7f1241b42e322a6451854ab7e4b4c146bf78a/lean/ComparatorChallenges/HyperbolicCones.lean}{Comparator HyperbolicCones.lean}.
\item \href{https://github.com/openai/math/blob/adc7f1241b42e322a6451854ab7e4b4c146bf78a/CONTENTS.md}{OpenAI catalogue entry}.
\end{itemize}
\end{document}
